% GENERATED FILE. Edit canonical lessons, then run npm run generate:books.
% canonical-source-hash: fnv1a64:2a57a3c85dcf603c
% canonical-lessons: ES-C375-una, ES-C375-labio, ES-C375-pulgar, ES-C375-beso, ES-C375-sobrino

\chapter{Five Things a Person Has}
\label{ch:la-una-y-el-labio}
\hlchaptermodality{375}

\begin{chapteropening}
\textbf{By the end of this chapter:} \emph{I can say la uña, el labio, el pulgar, el beso and el sobrino, and name five things that belong to a person.}

Complete ES-C375-sobrino: name five things that belong to a person, and say one thing about each.
\end{chapteropening}

\section[la uña]{la uña — a small claw}

\label{lesson:ES-C375-una}



\begin{quote}
Say \emph{the candle}, then \emph{the knot}. (\emph{La vela}; \emph{El nudo}.)
\end{quote}

\subsection*{What to know first}
\begin{itemize}
\raggedright
  \item el dedo — where it sits.
  \item la mano — the whole of which it is a corner.
\end{itemize}

\begin{sounds}
\begin{itemize}
\raggedright
  \item \emph{U-ña}, two beats with the weight on the first. The \emph{ñ} is one sound, the \emph{ny} in English \emph{canyon}, never two. $\to$ pronunciation reference
\end{itemize}
\end{sounds}

\begin{cousinweb}
\textbf{la uña} — \textquotedblleft{}fingernail.\textquotedblright{}

Latin \textbf{unguis} was a nail or a claw. \textbf{Ungula} was its diminutive: a small claw. Spanish inherited the diminutive, so \emph{uña} has a \textquotedblleft{}little\textquotedblright{} folded into it that no Spanish speaker hears any more.

English borrowed the same Latin diminutive much later and used it to name a whole branch of the animal kingdom. An \textbf{ungulate} is a hoofed animal — a horse, a deer, a cow — filed under a word that literally means \emph{small claw}, because a hoof is what a claw becomes when an animal walks on it.

One caution, because the resemblance invites the wrong conclusion. English \textbf{nail} is not descended from \emph{unguis}. The two are relatives at a depth that has to be reconstructed rather than read off the page, which is a different kind of claim from \emph{ungulate}, where you can watch the borrowing happen.
\end{cousinweb}

\begin{grammarlens}[title={What you've built}]
You can now name the hard part of a finger: \emph{la uña}.
\end{grammarlens}

\subsection*{Your turn}
Say these aloud:

\begin{itemize}
\raggedright
  \item \emph{la uña}
  \item \emph{las uñas}
  \item \emph{una uña rota}
\end{itemize}

\begin{tcolorbox}[breakable,colback=teal!4,colframe=teal!35!black,title={Before you move on}]
What did Latin \emph{ungula} mean? (A small claw — it was already a diminutive.) And what kind of animal is an ungulate? (A hoofed one, named from the same word.)
\end{tcolorbox}
\section[el labio]{el labio — a resemblance the handbooks have stopped explaining}

\label{lesson:ES-C375-labio}



\begin{quote}
Say \emph{the knot}, then \emph{the fingernail}. (\emph{El nudo}; \emph{La uña}.)
\end{quote}

\subsection*{What to know first}
\begin{itemize}
\raggedright
  \item la boca — what it closes.
  \item el diente — what sits behind it.
\end{itemize}

\begin{sounds}
\begin{itemize}
\raggedright
  \item \emph{LA-bio}, two beats with the weight on the first. The \emph{b} between vowels is soft — the lips come close without quite meeting. $\to$ pronunciation reference
\end{itemize}
\end{sounds}

\begin{cousinweb}
\textbf{el labio} — \textquotedblleft{}lip.\textquotedblright{}

Latin \textbf{labium} gave Spanish \emph{labio}, and English took the same word directly for its technical use: a \textbf{labial} sound is one made with the lips.

English \textbf{lip} is a different story. It is native Germanic, Old English \emph{lippa}, and it is not descended from \emph{labium}.

For a long time the handbooks put the two under one Indo-European root and left it there. That is no longer where the argument stands. De Vaan and Boutkan reject the Indo-European derivation for this whole group and suspect something stranger: that Latin and Germanic each \textbf{borrowed} the word from a language spoken in Europe before either of them arrived, and then lost.

So the resemblance between \emph{labio} and \emph{lip} is probably not an accident. It is probably also not inheritance. What it most likely records is a word that was already in Europe when the languages you can trace showed up — and a speaker nobody can name.
\end{cousinweb}

\begin{grammarlens}[title={What you've built}]
You can now name the edge of a mouth: \emph{el labio}.
\end{grammarlens}

\subsection*{Your turn}
Say these aloud:

\begin{itemize}
\raggedright
  \item \emph{el labio}
  \item \emph{los labios}
  \item \emph{el labio y el diente}
\end{itemize}

\begin{tcolorbox}[breakable,colback=teal!4,colframe=teal!35!black,title={Before you move on}]
Which English word is the direct borrowing of \emph{labium}? (\emph{Labial}.) And is English \emph{lip} descended from it? (No — Germanic, and the old shared-root story is now doubted.)
\end{tcolorbox}
\section[el pulgar]{el pulgar — a ruler before it was a word}

\label{lesson:ES-C375-pulgar}



\begin{quote}
Say \emph{the fingernail}, then \emph{the lip}. (\emph{La uña}; \emph{El labio}.)
\end{quote}

\subsection*{What to know first}
\begin{itemize}
\raggedright
  \item el dedo — the family it belongs to.
  \item la mano — where it sits, opposite the rest.
\end{itemize}

\begin{sounds}
\begin{itemize}
\raggedright
  \item \emph{pul-GAR}, two beats with the weight on the second. A word ending in \emph{-r} takes its stress on the last beat and needs no written accent. $\to$ pronunciation reference
\end{itemize}
\end{sounds}

\begin{cousinweb}
\textbf{el pulgar} — \textquotedblleft{}thumb.\textquotedblright{}

Latin \textbf{pollex} was the thumb. From it came the adjective \textbf{pollicāris} — and the Academy's dictionary glosses that adjective not as \textquotedblleft{}of the thumb\textquotedblright{} but as \textbf{\textquotedblleft{}an inch long.\textquotedblright{}}

That is the real find here. The thumb was a \textbf{measuring instrument} before it was a word for a body part, and Spanish still carries the proof: a \textbf{pulgada} is an inch. English does the same thing in reverse — \emph{inch} comes from Latin \emph{uncia}, a twelfth part, and \emph{thumb} stays out of it, but French \emph{pouce} means both thumb and inch to this day.

There is a story attached to \emph{pollex} that you will meet, and it should be set aside. It says the thumb is \textquotedblleft{}the strong one,\textquotedblright{} from the verb \emph{pollēre}, \textquotedblleft{}to be strong.\textquotedblright{}

Follow how that argument runs. It reads strength \textbf{out of} the proposed ancestor, then points at how strong a thumb is as \textbf{evidence for} that ancestor. The thumb's strength is a fact about thumbs; it proves nothing about the word. De Vaan is unconvinced, and \emph{pollex} is best left as disputed.
\end{cousinweb}

\begin{grammarlens}[title={What you've built}]
You can now name the odd one out on a hand: \emph{el pulgar}.
\end{grammarlens}

\subsection*{Your turn}
Say these aloud:

\begin{itemize}
\raggedright
  \item \emph{el pulgar}
  \item \emph{un pulgar}
  \item \emph{el pulgar y los dedos}
\end{itemize}

\begin{tcolorbox}[breakable,colback=teal!4,colframe=teal!35!black,title={Before you move on}]
What did Latin \emph{pollicāris} actually mean? (An inch long.) And which Spanish word for a measurement comes from it? (\emph{Pulgada}.)
\end{tcolorbox}
\section[el beso]{el beso — a word Latin borrowed too}

\label{lesson:ES-C375-beso}



\begin{quote}
Say \emph{the lip}, then \emph{the thumb}. (\emph{El labio}; \emph{El pulgar}.)
\end{quote}

\subsection*{What to know first}
\begin{itemize}
\raggedright
  \item el abrazo — what usually comes with it.
  \item la boca — what makes it.
\end{itemize}

\begin{sounds}
\begin{itemize}
\raggedright
  \item \emph{BE-so}, two beats with the weight on the first. The \emph{s} is clear and sharp, never buzzing into a \emph{z}. $\to$ pronunciation reference
\end{itemize}
\end{sounds}

\begin{cousinweb}
\textbf{el beso} — \textquotedblleft{}kiss.\textquotedblright{}

Latin \textbf{basium} gave Spanish \emph{beso}, French \textbf{baiser}, Italian \textbf{bacio}. Three Romance languages, one Latin noun, nothing surprising so far.

The surprise is in the Academy's own entry, which does not stop at Latin. It says \emph{basium} was itself \textbf{of Celtic origin} — so the word had already crossed a language boundary before Latin got hold of it, and every Romance kiss is a loan on top of a loan.

English has an old word \textbf{buss} for a kiss, and it is exactly the kind of resemblance that should make you slow down rather than speed up. \emph{Buss}, \emph{bacio}, Welsh and Gaelic \emph{bus} for a lip — they all sound alike because the sound of a kiss is a short one made with the lips, and several languages appear to have arrived at it separately.

That makes \emph{buss} a look-alike, not a cousin. A word can resemble another because of shared ancestry or because both are imitating the same noise, and only one of those is a family tie.
\end{cousinweb}

\begin{grammarlens}[title={What you've built}]
You can now name what goes with an embrace: \emph{un beso}.
\end{grammarlens}

\subsection*{Your turn}
Say these aloud:

\begin{itemize}
\raggedright
  \item \emph{el beso}
  \item \emph{un beso}
  \item \emph{un abrazo y un beso}
\end{itemize}

\begin{tcolorbox}[breakable,colback=teal!4,colframe=teal!35!black,title={Before you move on}]
Where did Latin get \emph{basium}? (From Celtic — Latin borrowed it too.) And is English \emph{buss} a relative? (Probably not; it imitates the sound.)
\end{tcolorbox}
\section[el sobrino]{el sobrino — a cousin that slid down a generation}

\label{lesson:ES-C375-sobrino}



\begin{quote}
Say \emph{the thumb}, then \emph{the kiss}. (\emph{El pulgar}; \emph{El beso}.)
\end{quote}

\subsection*{What to know first}
\begin{itemize}
\raggedright
  \item el primo — what the Latin word actually meant.
  \item el tío — the other end of the same relationship.
\end{itemize}

\begin{sounds}
\begin{itemize}
\raggedright
  \item \emph{so-BRI-no}, three beats with the weight on the middle one. The \emph{r} is a single tap, and \emph{br} runs together with no vowel between. $\to$ pronunciation reference
\end{itemize}
\end{sounds}

\begin{cousinweb}
\textbf{el sobrino} — \textquotedblleft{}nephew.\textquotedblright{}

Latin \textbf{sobrīnus} meant a \textbf{cousin} — on the mother's side. Underneath it is \textbf{soror}, \textquotedblleft{}sister\textquotedblright{}: \emph{sobrinus} is a worn-down \textbf{sororinus}, the child of your mother's sister.

Latin then built a second word on it. Prefixing \emph{com-} gave \textbf{consobrinus}, and that one named the \textbf{closer} relation, the first cousin, where bare \emph{sobrinus} drifted out to the more distant ones. Isidore of Seville uses \emph{consobrinus} for a first cousin and \emph{sobrinus} for a second.

Each of the two words then went its own way:

\begin{itemize}
\raggedright
  \item \textbf{consobrinus} travelled through Old French and became English \textbf{cousin};
  \item \textbf{sobrinus} stayed in Spanish and \textbf{moved down a generation}, from cousin to nephew.
\end{itemize}

So when an English speaker says \emph{cousin} and a Spanish speaker says \emph{sobrino}, they are reaching for the same Latin noun and landing on two different branches of the family. Spanish did not lose the cousin: \emph{el primo} took the job instead.
\end{cousinweb}

\begin{grammarlens}[title={What you've built}]
You can now name your sibling's child: \emph{mi sobrino}.
\end{grammarlens}

\subsection*{Your turn}
Say these aloud:

\begin{itemize}
\raggedright
  \item \emph{el sobrino}
  \item \emph{mi sobrino}
  \item \emph{mi sobrino y mi primo}
\end{itemize}

\begin{tcolorbox}[breakable,colback=teal!4,colframe=teal!35!black,title={Before you move on}]
What did Latin \emph{sobrīnus} mean? (A cousin, on the mother's side.) And which English word comes from its prefixed form \emph{consobrinus}? (\emph{Cousin}.)
\end{tcolorbox}
